\chapter{Meshing, Error and Convergence}
\label{ch:meshing-convergence}

A finite element result is only as good as its mesh. This chapter states
what convergence theory guarantees, what element quality demands, and how
to obtain trustworthy stresses --- the quantities engineers actually care
about, and the ones that converge slowest.

\section{A priori error estimates}

For an elliptic problem solved on a quasi-uniform mesh of size $h$ with
elements of complete polynomial degree $p$, and a sufficiently smooth
exact solution $\disp\in H^{p+1}$, the finite element error satisfies
\begin{align}
  \lVert \disp-\disp^h\rVert_{H^1} &\le C\,h^{p}\,|\disp|_{H^{p+1}}
    &&\text{(energy / }H^1\text{ norm)}, \label{eq:h1err}\\
  \lVert \disp-\disp^h\rVert_{L^2} &\le C\,h^{p+1}\,|\disp|_{H^{p+1}}
    &&\text{(}L^2\text{ norm, via Aubin--Nitsche duality)} . \label{eq:l2err}
\end{align}
Two consequences dominate practice. First, \textbf{stresses converge one
order slower than displacements}: stresses are first derivatives of
$\disp$, so they live in the $H^1$ (energy) norm with rate $p$, while
displacements themselves converge at rate $p+1$ in $L^2$. A displacement
that looks converged can still carry an under-resolved stress field.
Second, the rates in~\eqref{eq:h1err}--\eqref{eq:l2err} assume smoothness;
re-entrant corners, point loads and material interfaces introduce
singularities ($\disp\in H^{1+\alpha}$, $\alpha<1$) that \emph{reduce} the
achievable rate and produce mesh-dependent stress hot-spots.

\begin{fcwarn}[The corner-singularity trap]
At a sharp re-entrant corner the elastic stress is theoretically
\emph{infinite}; refining the mesh there makes the peak stress grow
without bound rather than converge. Do not chase such a number. Either
model the real fillet radius, or evaluate stress at a physically
meaningful distance from the singularity (structural-stress / hot-spot
methods).
\end{fcwarn}

\section{Consistency, stability and the patch test}

Convergence requires two ingredients: \textbf{consistency} (the element
can reproduce constant-strain and rigid-body states --- completeness) and
\textbf{stability}. The practical check on consistency is the
\textbf{patch test} (Irons): subject a patch of (possibly distorted)
elements to boundary displacements corresponding to a constant strain
field; a passing element reproduces that constant stress exactly, to
machine precision. Passing the patch test certifies that the element will
converge under refinement, and it is the standard validation for
nonconforming and reduced-integration elements.

\section{$h$-, $p$- and $hp$-refinement}

\begin{description}
  \item[$h$-refinement] shrinks $h$ at fixed $p$: algebraic convergence
        $\mathcal O(h^p)$. The bread-and-butter of automatic meshers.
  \item[$p$-refinement] raises the polynomial degree at fixed mesh:
        \emph{exponential} convergence for smooth solutions, algebraic
        near singularities.
  \item[$hp$-refinement] combines them --- small, low-order elements at
        singularities, large high-order elements in smooth regions ---
        recovering exponential rates even with singularities.
\end{description}
In FreeCAD/CalculiX practice you have a pragmatic version of this choice:
select second-order elements (Tet10, effectively $p=2$) and then refine
$h$ locally where gradients are steep.

\section{Element quality}

Distortion degrades both accuracy and conditioning. The mesh metrics to
watch:
\begin{itemize}
  \item \textbf{Aspect ratio} --- ratio of longest to shortest element
        dimension; keep modest (rules of thumb: $<5$--$10$ in stressed
        regions).
  \item \textbf{Skew / interior angles} --- angles far from the ideal
        (60\textdegree{} for triangles, 90\textdegree{} for quads)
        increase error; avoid angles near 0\textdegree{} or 180\textdegree.
  \item \textbf{Taper and warping} --- for quad/hex faces.
  \item \textbf{Jacobian} --- $\det\Jmat$ must stay positive at every
        integration point; a negative Jacobian is an inverted element and
        a hard error.
\end{itemize}
A locally refined mesh of well-shaped elements beats a uniformly fine mesh
of distorted ones, for both cost and accuracy.

\section{Stress recovery and error estimation}

Stresses computed from $\svec=\Dmat\Bmat\dof$ are \emph{not} equally
accurate everywhere in an element. They are most accurate at special
interior points --- the \textbf{Barlow / reduced-Gauss points} --- where
they exhibit \textbf{superconvergence} (one order better than at nodes).
Practical post-processors exploit this:

\begin{description}
  \item[Superconvergent Patch Recovery (SPR).] The Zienkiewicz--Zhu
        procedure~\citep{zienkiewicz1992spr} least-squares-fits a smooth
        polynomial stress field to the superconvergent sampling points
        over a patch of elements, then evaluates it at nodes --- yielding
        markedly improved, continuous stresses.
  \item[The $Z^2$ (ZZ) error estimator.] The difference between the raw
        (discontinuous) element stresses and the recovered (smooth) field
        estimates the local discretisation error a posteriori, and drives
        \textbf{adaptive remeshing}: refine where the estimated error is
        large, coarsen where it is small.
\end{description}

\begin{fctip}[A convergence workflow you can trust]
\begin{enumerate}
  \item Solve on a coarse mesh; note the quantity of interest (peak von
        Mises stress, a key displacement).
  \item Refine (globally, or locally where SPR/ZZ or raw stress
        discontinuities flag error) and re-solve.
  \item Repeat until the quantity of interest changes by less than your
        tolerance (say, a few percent) between refinements --- that is
        \emph{mesh convergence}.
  \item Watch for non-convergence at singularities
        (\S\ref{sec:corner}); a monotonically growing peak stress is a
        red flag, not a result.
\end{enumerate}
Never report a single-mesh stress without having demonstrated its mesh
insensitivity.
\end{fctip}
\label{sec:corner}

This closes the theory. Part~III now maps every construct above ---
elements, materials, constraints, solvers, stress recovery --- onto the
concrete objects and dialogs of FreeCAD's FEM Workbench.
